% Seminar 2 - Pagina 6. Fragment pentru % Seminar 2 - Pagina 6. Fragment pentru % Seminar 2 - Pagina 6. Fragment pentru % Seminar 2 - Pagina 6. Fragment pentru \input{seminars/pages/seminar2-pagina6-ro}.
% Fara preambul si fara intrerupere fortata de pagina.
% NOTE EDITORIALE:
% - Continuarea exercitiului 7 incepe la mijlocul demonstratiei de pe pagina 5.
% - Unele relatii intermediare sunt taiate / suprascrise; doar relatiile
%   lizibile si netaiate sunt incluse in transcrierea principala.
% - In argumentul prin sir descrescator, indicii u_1,u_2,u_3 sunt cititi
%   conform scrisului, dar cuantificarea este prescurtata in manuscris.
% - La exercitiul 8, simbolul Hom(A,B) este utilizat pentru toate
%   aplicatiile A -> B; nu este vorba explicit despre morfisme algebrice.
% - Exercitiul 9 afirma inexistenta unei surjectii X -> P(X), dar nu
%   contine demonstratie pe aceasta pagina.
% - La exercitiul 10 este demonstrata doar injectivitatea lui phi;
%   nu este adaugata o dovada a surjectivitatii.

\subsection*{Seminar 2 - Pagina 6}

\noindent\underline{Sol. 7 (continuare).}
\[
\Longrightarrow f(n)\geq k+1.
\]
\noindent Deci
\[
 f(k)\geq k\qquad (\forall k\in\mathbb{N}^{*}).
\]

\noindent Presupunem prin reducere la absurd că există
$k\in\mathbb{N}^{*}$ a.î.
\[
 f(k)>k.
\]
\noindent Arătăm că, dacă $n>k$, atunci $f(n-1)>k$.
\[
 n>k\ \Longrightarrow\ n\geq k+1\ \Longrightarrow\ n-1\geq k
 \ \Longrightarrow\ f(n-1)\geq k.
\]
\noindent Presupunem $f(n-1)=k$. Atunci
\[
 k=f(n-1)\geq n-1\geq k
 \quad\Longrightarrow\quad n-1=k
 \quad\Longrightarrow\quad f(k)=k,\quad\text{fals}.
\]

\noindent Fie $u_1>k$. Atunci
\[
 f(u_1)>f(u_2),\qquad u_2=f(u_1-1)>k.
\]
\noindent În mod similar,
\[
 f(u_2)>f(u_3),\qquad u_3=f(u_2-1)>k.
\]
\noindent Continuând la infinit, obţinem un şir infinit strict
 descrescător de numere naturale.

\bigskip
\noindent\fbox{8}\quad
\[
 \operatorname{Hom}(A,B)\overset{\mathrm{def}}{=}
 \{f\mid f:A\longrightarrow B\}.
\]
\noindent $\mathcal{P}(X)$ în bijecţie cu
\[
 \operatorname{Hom}(X,\{0,1\})
   =\{\varphi\mid\varphi:X\longrightarrow\{0,1\}\}.
\]

\medskip
\noindent\fbox{9}\quad Nu există aplicaţie surjectivă
\[
 f:X\longrightarrow\mathcal{P}(X).
\]

\medskip
\noindent\fbox{10}\quad
\[
 \varphi:\mathcal{P}(X)\longrightarrow
 \operatorname{Hom}(X,\{0,1\}).
\]
\noindent Dacă $A\subseteq X$, atunci
$\varphi(A)=\varphi_A:X\longrightarrow\{0,1\}$, unde
\[
 \varphi_A(x)=
 \begin{cases}
  1,&x\in A,\\
  0,&x\notin A.
 \end{cases}
\]
\noindent $\varphi$ este injectivă: dacă
\[
 \varphi(A)=\varphi(B),\qquad \varphi_A=\varphi_B,
\]
\noindent atunci
\[
 x\in A\ \Longleftrightarrow\ \varphi_A(x)=1
 \ \Longleftrightarrow\ \varphi_B(x)=1
 \ \Longleftrightarrow\ x\in B
 \quad\Longrightarrow\quad A=B.
\]
.
% Fara preambul si fara intrerupere fortata de pagina.
% NOTE EDITORIALE:
% - Continuarea exercitiului 7 incepe la mijlocul demonstratiei de pe pagina 5.
% - Unele relatii intermediare sunt taiate / suprascrise; doar relatiile
%   lizibile si netaiate sunt incluse in transcrierea principala.
% - In argumentul prin sir descrescator, indicii u_1,u_2,u_3 sunt cititi
%   conform scrisului, dar cuantificarea este prescurtata in manuscris.
% - La exercitiul 8, simbolul Hom(A,B) este utilizat pentru toate
%   aplicatiile A -> B; nu este vorba explicit despre morfisme algebrice.
% - Exercitiul 9 afirma inexistenta unei surjectii X -> P(X), dar nu
%   contine demonstratie pe aceasta pagina.
% - La exercitiul 10 este demonstrata doar injectivitatea lui phi;
%   nu este adaugata o dovada a surjectivitatii.

\subsection*{Seminar 2 - Pagina 6}

\noindent\underline{Sol. 7 (continuare).}
\[
\Longrightarrow f(n)\geq k+1.
\]
\noindent Deci
\[
 f(k)\geq k\qquad (\forall k\in\mathbb{N}^{*}).
\]

\noindent Presupunem prin reducere la absurd că există
$k\in\mathbb{N}^{*}$ a.î.
\[
 f(k)>k.
\]
\noindent Arătăm că, dacă $n>k$, atunci $f(n-1)>k$.
\[
 n>k\ \Longrightarrow\ n\geq k+1\ \Longrightarrow\ n-1\geq k
 \ \Longrightarrow\ f(n-1)\geq k.
\]
\noindent Presupunem $f(n-1)=k$. Atunci
\[
 k=f(n-1)\geq n-1\geq k
 \quad\Longrightarrow\quad n-1=k
 \quad\Longrightarrow\quad f(k)=k,\quad\text{fals}.
\]

\noindent Fie $u_1>k$. Atunci
\[
 f(u_1)>f(u_2),\qquad u_2=f(u_1-1)>k.
\]
\noindent În mod similar,
\[
 f(u_2)>f(u_3),\qquad u_3=f(u_2-1)>k.
\]
\noindent Continuând la infinit, obţinem un şir infinit strict
 descrescător de numere naturale.

\bigskip
\noindent\fbox{8}\quad
\[
 \operatorname{Hom}(A,B)\overset{\mathrm{def}}{=}
 \{f\mid f:A\longrightarrow B\}.
\]
\noindent $\mathcal{P}(X)$ în bijecţie cu
\[
 \operatorname{Hom}(X,\{0,1\})
   =\{\varphi\mid\varphi:X\longrightarrow\{0,1\}\}.
\]

\medskip
\noindent\fbox{9}\quad Nu există aplicaţie surjectivă
\[
 f:X\longrightarrow\mathcal{P}(X).
\]

\medskip
\noindent\fbox{10}\quad
\[
 \varphi:\mathcal{P}(X)\longrightarrow
 \operatorname{Hom}(X,\{0,1\}).
\]
\noindent Dacă $A\subseteq X$, atunci
$\varphi(A)=\varphi_A:X\longrightarrow\{0,1\}$, unde
\[
 \varphi_A(x)=
 \begin{cases}
  1,&x\in A,\\
  0,&x\notin A.
 \end{cases}
\]
\noindent $\varphi$ este injectivă: dacă
\[
 \varphi(A)=\varphi(B),\qquad \varphi_A=\varphi_B,
\]
\noindent atunci
\[
 x\in A\ \Longleftrightarrow\ \varphi_A(x)=1
 \ \Longleftrightarrow\ \varphi_B(x)=1
 \ \Longleftrightarrow\ x\in B
 \quad\Longrightarrow\quad A=B.
\]
.
% Fara preambul si fara intrerupere fortata de pagina.
% NOTE EDITORIALE:
% - Continuarea exercitiului 7 incepe la mijlocul demonstratiei de pe pagina 5.
% - Unele relatii intermediare sunt taiate / suprascrise; doar relatiile
%   lizibile si netaiate sunt incluse in transcrierea principala.
% - In argumentul prin sir descrescator, indicii u_1,u_2,u_3 sunt cititi
%   conform scrisului, dar cuantificarea este prescurtata in manuscris.
% - La exercitiul 8, simbolul Hom(A,B) este utilizat pentru toate
%   aplicatiile A -> B; nu este vorba explicit despre morfisme algebrice.
% - Exercitiul 9 afirma inexistenta unei surjectii X -> P(X), dar nu
%   contine demonstratie pe aceasta pagina.
% - La exercitiul 10 este demonstrata doar injectivitatea lui phi;
%   nu este adaugata o dovada a surjectivitatii.

\subsection*{Seminar 2 - Pagina 6}

\noindent\underline{Sol. 7 (continuare).}
\[
\Longrightarrow f(n)\geq k+1.
\]
\noindent Deci
\[
 f(k)\geq k\qquad (\forall k\in\mathbb{N}^{*}).
\]

\noindent Presupunem prin reducere la absurd că există
$k\in\mathbb{N}^{*}$ a.î.
\[
 f(k)>k.
\]
\noindent Arătăm că, dacă $n>k$, atunci $f(n-1)>k$.
\[
 n>k\ \Longrightarrow\ n\geq k+1\ \Longrightarrow\ n-1\geq k
 \ \Longrightarrow\ f(n-1)\geq k.
\]
\noindent Presupunem $f(n-1)=k$. Atunci
\[
 k=f(n-1)\geq n-1\geq k
 \quad\Longrightarrow\quad n-1=k
 \quad\Longrightarrow\quad f(k)=k,\quad\text{fals}.
\]

\noindent Fie $u_1>k$. Atunci
\[
 f(u_1)>f(u_2),\qquad u_2=f(u_1-1)>k.
\]
\noindent În mod similar,
\[
 f(u_2)>f(u_3),\qquad u_3=f(u_2-1)>k.
\]
\noindent Continuând la infinit, obţinem un şir infinit strict
 descrescător de numere naturale.

\bigskip
\noindent\fbox{8}\quad
\[
 \operatorname{Hom}(A,B)\overset{\mathrm{def}}{=}
 \{f\mid f:A\longrightarrow B\}.
\]
\noindent $\mathcal{P}(X)$ în bijecţie cu
\[
 \operatorname{Hom}(X,\{0,1\})
   =\{\varphi\mid\varphi:X\longrightarrow\{0,1\}\}.
\]

\medskip
\noindent\fbox{9}\quad Nu există aplicaţie surjectivă
\[
 f:X\longrightarrow\mathcal{P}(X).
\]

\medskip
\noindent\fbox{10}\quad
\[
 \varphi:\mathcal{P}(X)\longrightarrow
 \operatorname{Hom}(X,\{0,1\}).
\]
\noindent Dacă $A\subseteq X$, atunci
$\varphi(A)=\varphi_A:X\longrightarrow\{0,1\}$, unde
\[
 \varphi_A(x)=
 \begin{cases}
  1,&x\in A,\\
  0,&x\notin A.
 \end{cases}
\]
\noindent $\varphi$ este injectivă: dacă
\[
 \varphi(A)=\varphi(B),\qquad \varphi_A=\varphi_B,
\]
\noindent atunci
\[
 x\in A\ \Longleftrightarrow\ \varphi_A(x)=1
 \ \Longleftrightarrow\ \varphi_B(x)=1
 \ \Longleftrightarrow\ x\in B
 \quad\Longrightarrow\quad A=B.
\]
.
% Fara preambul si fara intrerupere fortata de pagina.
% NOTE EDITORIALE:
% - Continuarea exercitiului 7 incepe la mijlocul demonstratiei de pe pagina 5.
% - Unele relatii intermediare sunt taiate / suprascrise; doar relatiile
%   lizibile si netaiate sunt incluse in transcrierea principala.
% - In argumentul prin sir descrescator, indicii u_1,u_2,u_3 sunt cititi
%   conform scrisului, dar cuantificarea este prescurtata in manuscris.
% - La exercitiul 8, simbolul Hom(A,B) este utilizat pentru toate
%   aplicatiile A -> B; nu este vorba explicit despre morfisme algebrice.
% - Exercitiul 9 afirma inexistenta unei surjectii X -> P(X), dar nu
%   contine demonstratie pe aceasta pagina.
% - La exercitiul 10 este demonstrata doar injectivitatea lui phi;
%   nu este adaugata o dovada a surjectivitatii.

\subsection*{Seminar 2 - Pagina 6}

\noindent\underline{Sol. 7 (continuare).}
\[
\Longrightarrow f(n)\geq k+1.
\]
\noindent Deci
\[
 f(k)\geq k\qquad (\forall k\in\mathbb{N}^{*}).
\]

\noindent Presupunem prin reducere la absurd că există
$k\in\mathbb{N}^{*}$ a.î.
\[
 f(k)>k.
\]
\noindent Arătăm că, dacă $n>k$, atunci $f(n-1)>k$.
\[
 n>k\ \Longrightarrow\ n\geq k+1\ \Longrightarrow\ n-1\geq k
 \ \Longrightarrow\ f(n-1)\geq k.
\]
\noindent Presupunem $f(n-1)=k$. Atunci
\[
 k=f(n-1)\geq n-1\geq k
 \quad\Longrightarrow\quad n-1=k
 \quad\Longrightarrow\quad f(k)=k,\quad\text{fals}.
\]

\noindent Fie $u_1>k$. Atunci
\[
 f(u_1)>f(u_2),\qquad u_2=f(u_1-1)>k.
\]
\noindent În mod similar,
\[
 f(u_2)>f(u_3),\qquad u_3=f(u_2-1)>k.
\]
\noindent Continuând la infinit, obţinem un şir infinit strict
 descrescător de numere naturale.

\bigskip
\noindent\fbox{8}\quad
\[
 \operatorname{Hom}(A,B)\overset{\mathrm{def}}{=}
 \{f\mid f:A\longrightarrow B\}.
\]
\noindent $\mathcal{P}(X)$ în bijecţie cu
\[
 \operatorname{Hom}(X,\{0,1\})
   =\{\varphi\mid\varphi:X\longrightarrow\{0,1\}\}.
\]

\medskip
\noindent\fbox{9}\quad Nu există aplicaţie surjectivă
\[
 f:X\longrightarrow\mathcal{P}(X).
\]

\medskip
\noindent\fbox{10}\quad
\[
 \varphi:\mathcal{P}(X)\longrightarrow
 \operatorname{Hom}(X,\{0,1\}).
\]
\noindent Dacă $A\subseteq X$, atunci
$\varphi(A)=\varphi_A:X\longrightarrow\{0,1\}$, unde
\[
 \varphi_A(x)=
 \begin{cases}
  1,&x\in A,\\
  0,&x\notin A.
 \end{cases}
\]
\noindent $\varphi$ este injectivă: dacă
\[
 \varphi(A)=\varphi(B),\qquad \varphi_A=\varphi_B,
\]
\noindent atunci
\[
 x\in A\ \Longleftrightarrow\ \varphi_A(x)=1
 \ \Longleftrightarrow\ \varphi_B(x)=1
 \ \Longleftrightarrow\ x\in B
 \quad\Longrightarrow\quad A=B.
\]
